\chapter{FGA representability of the Picard scheme}
\label{chap:Picard_FGAPicRepresentability}
% archon:covers AlgebraicJacobian/Picard/FGAPicRepresentability.lean

\section{Setup and motivation}
\label{sec:fga_pic_setup}

Throughout this chapter, fix a base field \(k\) and a smooth proper geometrically
integral curve \(C\) over \(k\) (the curve carrying the Jacobian). The relative
Picard presheaf \(\Pic^\sharp_{C/k} : (\Sch/k)^{op} \to \Set\) of
\cref{chap:Picard_LineBundlePullback} has already been packaged as a
set-valued functor and refined to its abelian-group-valued
\'etale-sheafified version \(\Pic^\sharp_{(C/k)\et}\) in
\cref{chap:Picard_RelPicFunctor}. What remains is to produce a scheme
representing it: the \emph{Picard scheme} \(\Pic_{C/k}\). The present chapter
packages that assembly, together with the abelian-group-scheme refinement.

\medskip

\textbf{Two routes to representability.} The Picard scheme of a curve can be
constructed in two ways, and this blueprint records both, because they need
different inputs and only one of them is the route under construction here.

The \emph{quotient route} is Grothendieck's, as presented by Kleiman \S 4: the
Abel map \(\mathrm{Div}_{C/k} \to \Pic^\sharp_{C/k}\) exhibits the Picard functor
as a quotient of a quasi-projective scheme by a smooth proper equivalence
relation, and Altman--Kleiman descent
(\cref{lem:smooth_proper_quotient}) represents the quotient. Its engine is the
Hilbert scheme \(\Hilb_{C/k} = \Quot_{\mathcal{O}_C/C/k}\) of
\cref{chap:Picard_QuotScheme}, whose open subscheme \(\mathrm{Div}_{C/k}\) of
relative effective divisors carries the source of the Abel map. This is the route
of \cref{thm:fga_pic_representability} and its proof below.

The \emph{Milne--Koll\'ar route} builds \(\Pic^r_{C/k}\) directly, over a
separably closed field, from the open loci
\[
  C^\Sigma = \{\,t : h^0(\mathcal{O}(W - \Sigma)_t) = 1\,\}
\]
cut out inside a degree-\(r\) divisor scheme by a tuple \(\Sigma\) of rational
points, glues them, and descends the result to \(k\) by a finite Galois quotient
before assembling the degrees by a coproduct. Its inputs are a rigidified
pushforward engine, uniform \(H^1\) vanishing, and Galois descent of rigidified
classes --- not a Hilbert or Quot scheme.

The decisive difference is the quotient step. Altman--Kleiman descent needs
quasi-projectivity of the scheme being quotiented, which is not expressible in
the formalisation at present and without which the lemma is false
(\cref{rem:smooth_proper_quotient_hypothesis}); the Milne--Koll\'ar route takes
only finite Galois quotients under an orbit-in-affine hypothesis, and so avoids
that obstruction entirely. The Lean development therefore pursues the
Milne--Koll\'ar route. The quotient-route material below is retained as the
mathematics it is --- the divisor and Hilbert-scheme substrate it rests on is
consumed by both routes --- and is not the path currently being formalised.

\medskip

\textbf{Comparison with the sheafified Picard functor.} Kleiman \S 4 represents
the \'etale-sheafified functor \(\Pic_{(C/k)\et}\). This chapter instead
works throughout under the standing hypothesis that \emph{\(C\) has a
\(k\)-rational point} (\cref{def:has_rational_point}) and states
representability against the \emph{plain} relative Picard functor
\(\Pic^\sharp_{C/k}\) of \cref{def:rel_pic_sharp}. This is justified by
Kleiman \S 2, Thm.~2.5 (\texttt{th:comp}): if \(f : X \to S\) has a section
and \(\mathcal O_{S'} = f'_* \mathcal O_{X'}\) holds for all base changes
(automatic for our proper geometrically integral curve, by cohomology and
base change), then the comparison maps
\[
  \Pic_{X/S}(T) \longrightarrow \Pic_{(X/S)\et}(T) \longrightarrow
  \Pic_{(X/S)\fppf}(T)
\]
are bijective for every \(T\). Hence under the rational-point hypothesis the
plain relative functor \emph{is} the \'etale sheaf, so the rational point
allows us to work directly with \(\Pic^\sharp_{C/k}\).

The hypothesis is not a convenience: without a section the plain functor is not
even a Zariski sheaf in general, so a representability statement for
\(\Pic^\sharp_{C/k}\) with no rational point would be false. Two honest
formulations are therefore available, and they prove different theorems.

\begin{enumerate}
  \item Represent \(\Pic^\sharp_{C/k}\) and carry \(C(k) \neq \emptyset\) as a
    hypothesis. This is the formulation of
    \cref{thm:fga_pic_representability} and of every representability assertion
    in this chapter. It is strictly weaker than a statement about all smooth
    proper geometrically integral curves, since such a curve need not have a
    rational point --- a conic without rational points, or a genus-\(2\) curve
    over \(\mathbb{Q}\) with no rational point, satisfies every other
    hypothesis.
  \item Represent the \'etale sheafification \(\Pic_{(C/k)\et}\) and drop the
    rational point. This is Kleiman's own formulation, and needs no section
    precisely because sheafifying supplies \'etale-locally what the section
    would have supplied globally. It requires developing the Picard functor on
    the \'etale site.
\end{enumerate}

Which formulation the project should target is a decision about what theorem is
being claimed, not a limitation of the platform. The two are recorded here
side by side; the chapter as written realises the first.

Every representability assertion of this chapter accordingly carries the
rational-point hypothesis. A pointed curve \((C,P)\), as in
\cref{chap:AbelJacobi}, supplies such a section.

% SOURCE: [Kleiman], ``The Picard scheme'', \S 4 (FGA Explained Ch.~9 \S 9.4);
% arXiv:math/0504020, p.~26. Read from
% references/kleiman-picard-src/kleiman-picard.tex.

\section{Line bundles and the Quot / Hilbert scheme}
\label{sec:fga_pic_line_bundle_quot}

The bridge from line bundles on \(C \times_k T\) to the Quot scheme is the
Abel map: a relative effective divisor \(D \subset C \times_k T\) over \(T\)
(a \(T\)-flat closed subscheme whose ideal is invertible) determines a closed
subscheme of \(C \times_k T\) flat over \(T\), hence a \(T\)-point of the Hilbert
functor \(\Hilb_{C/k}(T)\), which is the special case \(\Quot_{\mathcal{O}_C/C/k}\)
of the Quot functor of \cref{chap:Picard_QuotScheme}; and it determines an
invertible sheaf \(\mathcal{O}_{C \times_k T}(D)\) on the relative curve,
representing a class in \(\Pic^\sharp_{C/k}(T)\).

\begin{definition}[Abel-map witness]\leanok
  \label{def:abel_map_witness}
  \lean{AlgebraicGeometry.Scheme.PicScheme.abelMapWitness}
  \uses{def:div_functor_carrier, def:pic_sharp_carrier}
  Let \(C/k\) be a smooth proper curve. Define a
  natural transformation
  \[
    A^{\mathrm{raw}}_{C/k} \colon \mathrm{Div}_{C/k}
      \longrightarrow \Pic^\sharp_{C/k}
  \]
  as the composite
  \(
    \mathrm{abelKernelNatTrans}(C)
    \mathbin{;} \mathrm{picNeg}(C)
  \).
\end{definition}

\begin{proof}\leanok
  The first transformation sends the class represented by
  \(\langle \mathcal F,q\rangle\) to \([\ker q]\). It is well defined because
  equivalent divisor families induce isomorphic kernels. It is natural under
  arbitrary base change because the kernel--pullback comparison is an
  isomorphism for an epimorphic quotient with finitely presented, base-flat
  target and invertible kernel. The second transformation is pointwise
  inversion in the group-valued presheaf \(\Pic^\sharp_{C/k}\). Their composite
  is therefore natural and sends
  \(\langle \mathcal F,q\rangle\) to \(-[\ker q]\).
\end{proof}

\begin{lemma}[Line bundle / Quot correspondence via the Abel map]\leanok
  \label{lem:line_bundle_quot_correspondence}
  \lean{AlgebraicGeometry.Scheme.PicScheme.abelMap_app_mk}
  \uses{def:abel_map_witness, def:inst_has_abel_map, def:has_abel_map,
    def:div_functor_carrier, def:pic_sharp_carrier}
  % SOURCE: [Kleiman], ``The Picard scheme'', \S 3, Def.~dfn:Abel (the Abel
  % map), and Thm.~th:repDiv (representability of relative effective
  % divisors by an open subscheme of \(\IHilb_{X/S}\)); arXiv:math/0504020,
  % pp.~25, 23 (read from
  % references/kleiman-picard-src/kleiman-picard.tex, L1837-L1841 + L2135-L2145).
  % Source identifiers: \begin{thm}\label{th:repDiv}\end{thm}
  % and \begin{dfn}\label{dfn:Abel}\end{dfn}.
  \textit{Source: [Kleiman], ``The Picard scheme'', \S 3,
  Thm.~3.7 + Def.~4.6 (cf.\ [Nitsure], ``Construction of Hilbert
  and Quot Schemes'', \S 1, where \(\Hilb_{X/S} = \Quot_{\mathcal{O}_X/X/S}\)).}
  Let \(C/k\) be a smooth proper geometrically integral curve. There is a
  natural transformation of presheaves on \((\Sch/k)^{op}\)
  \[
    A_{C/k} = \mathrm{abelMap} \;\colon\; \mathrm{Div}_{C/k}
      \;\longrightarrow\; \Pic^\sharp_{C/k}
  \]
  (\(\mathrm{Div}_{C/k}\) the relative-divisor functor of
  \cref{def:div_functor_carrier}, \(\Pic^\sharp_{C/k}\) the relative
  Picard functor of \cref{def:pic_sharp_carrier}), with the following
  \emph{defining property}: at a \(T\)-point represented by a family
  \(\langle \mathcal F, q\rangle\) (\cref{def:div_family}: a finitely
  presented, \(T\)-flat quotient
  \(q \colon \mathcal O_{C\times_kT}\twoheadrightarrow\mathcal F\)
  whose kernel is an invertible ideal sheaf cutting out a relative effective
  divisor \(D\subseteq C\times_kT\)),
  \[
    A_{C/k}(T)\bigl(\langle \mathcal F, q\rangle\bigr)
      \;=\; -[\ker q] \;=\; [(\ker q)^{-1}]
      \;=\; \bigl[\mathcal{O}_{C \times_k T}(D)\bigr]
      \;\in\; \Pic^\sharp_{C/k}(T),
  \]
  the class of the associated invertible sheaf \(\mathcal{O}_{C \times_k
  T}(D)\), i.e.\ the additive inverse of the class of the ideal sheaf
  \(\ker q\) (Kleiman \S 3, Def.~dfn:Abel). This transformation is the
  required witness for \(\mathrm{HasAbelMap}(C)\)
  (\cref{def:has_abel_map}).
\end{lemma}

\begin{proof}
  The instance of \cref{def:inst_has_abel_map} identifies
  \(A_{C/k}\) with the witness of \cref{def:abel_map_witness}. Evaluating its
  two factors at \(\langle \mathcal F,q\rangle\) first gives the kernel class
  \([\ker q]\) and then its inverse. Since the kernel is the ideal sheaf of
  \(D\), its inverse is \(\mathcal O_{C\times_kT}(D)\), which gives the stated
  formula.
\end{proof}

\begin{remark}
[Representability of \(\mathrm{Div}_{C/k}\)]
  \label{rem:div_functor_representability}
  Kleiman~\S 3, Thm.~th:repDiv, shows that \(\mathrm{Div}_{C/k}\) is
  representable by an open subscheme \(\mathrm{Div}_{C/k} \subseteq
  \Hilb_{C/k}\) of the Hilbert scheme (\(\Hilb_{C/k} =
  \Quot_{\mathcal{O}_C/C/k}\), \cref{chap:Picard_QuotScheme}): letting
  \(W \subseteq C \times_k \Hilb_{C/k}\) be the universal closed
  subscheme and \(q \colon W \to \Hilb_{C/k}\) the projection, the locus
  \(V \subseteq W\) where \(W\) is locally an effective Cartier divisor
  is open, so \(Z := q(W \setminus V)\) is closed (\(q\) proper) and
  \(\mathrm{Div}_{C/k} := \Hilb_{C/k} \setminus Z\) is the open
  subscheme representing relative effective divisors. This is
  the relative-divisor functor appearing in
  \cref{def:div_functor_carrier}.
\end{remark}

\section{The FGA representability theorem}
\label{sec:fga_pic_representability}

We now state and prove the specialisation of Grothendieck's existence theorem
for the Picard scheme, in the specialisation needed for the Jacobian:
\(S = \Spec k\), \(X = C\) a smooth proper geometrically integral
curve. The general theorem (Kleiman \S 4, Thm.~4.8) covers any
\(X/S\) projective Zariski-locally over \(S\), flat, with integral geometric
fibers; we record the general statement and then specialise.

\begin{theorem}[FGA representability of the Picard scheme]
  \label{thm:fga_pic_representability}
  \lean{AlgebraicGeometry.Scheme.PicScheme.representable}
  \uses{lem:line_bundle_quot_correspondence, lem:smooth_proper_quotient,
        thm:quot_representable, def:rel_pic_sharp, def:has_rational_point,
        cor:flattening_stratification_curve, def:pic_sharp_representable,
        def:inst_pic_sharp_representable, lem:cech_flat_base_change}
  % SOURCE: [Kleiman], ``The Picard scheme'', \S 4, Main Thm.~th:main,
  % and Cor.~cor:algsch (specialisation to \(S = \Spec k\), \(X\) complete);
  % arXiv:math/0504020, pp.~26, 35 (read from
  % references/kleiman-picard-src/kleiman-picard.tex, L2155-L2166 + L2686-L2690).
  % Source identifiers: \begin{thm}\label{th:main}\end{thm}
  % and \begin{sbscor}\label{cor:algsch}\end{sbscor}.
  \textit{Source: [Kleiman], ``The Picard scheme'', \S 4, Main Thm.~4.8
  + Cor.~4.18.3.}
  Let \(k\) be a field and \(C\) a smooth proper geometrically integral curve
  over \(k\) \emph{with a \(k\)-rational point}
  (\cref{def:has_rational_point}). Then the relative Picard functor
  \(\Pic^\sharp_{C/k}\) of \cref{def:rel_pic_sharp} --- which under the
  rational-point hypothesis coincides with its \'etale sheafification
  \(\Pic_{(C/k)\et}\), by Kleiman \S 2, Thm.~2.5 --- is representable by a
  \(k\)-scheme \(\Pic_{C/k}\), separated and locally of finite type over
  \(k\), which is a disjoint union of open quasi-projective
  \(k\)-subschemes, each an increasing union of open quasi-projective
  subschemes. The scheme \(\Pic_{C/k}\) is called the \emph{Picard scheme}
  of \(C/k\).
\end{theorem}

\begin{proof}
  
  We follow Kleiman~\S 4, Thm.~4.8, specialised to \(S = \Spec k\) and
  \(X = C\). The proof proceeds in four conceptual steps, preceded by a
  comparison step; the underlying algebraic engine is the Quot scheme of
  \cref{chap:Picard_QuotScheme} and the relative Picard functor of
  \cref{chap:Picard_RelPicFunctor}, together with the line-bundle / Quot
  correspondence of \cref{lem:line_bundle_quot_correspondence}.

  \emph{Step 0: reduce to the \'etale sheaf via the rational point.}
  By Kleiman \S 2, Thm.~2.5, since \(C \to \Spec k\) has a section
  (\cref{def:has_rational_point}) and \(\mathcal O_T = \pi_{T*}
  \mathcal O_{C \times_k T}\) holds for every \(k\)-scheme \(T\) (cohomology
  and base change for the proper geometrically integral fibers of
  \(C\times_k T/T\)), the comparison map
  \(\Pic^\sharp_{C/k}(T) \to \Pic_{(C/k)\et}(T)\) is bijective for all
  \(T\). It therefore suffices to represent the \'etale sheaf
  \(\Pic_{(C/k)\et}\), which is what Kleiman's theorem produces; the same
  scheme then represents \(\Pic^\sharp_{C/k}\).

  \emph{Step 1: stratify by Hilbert polynomial.}
  For each polynomial \(\phi \in \mathbb{Q}[n]\), let \(P^\phi \subseteq
  \Pic^\sharp_{(C/k)\et}\) be the \'etale subsheaf whose \(T\)-points are
  represented by invertible sheaves \(\mathcal{L}\) on \(C \times_k T\) with
  \(\chi(C_t, \mathcal{L}_t^{-1}(n)) = \phi(n)\) for all \(t \in T\) (the
  Hilbert polynomial relative to the chosen polarisation
  \(\mathcal{O}_C(1)\) on the projective curve \(C\)). The \(P^\phi\) form a
  disjoint open cover of \(\Pic^\sharp_{(C/k)\et}\) (openness via Kleiman's
  flat-base-change argument, \cref{lem:cech_flat_base_change} / Stacks Tag 02KH;
  cf.\ EGA III 7.9.11). It suffices to
  represent each \(P^\phi\) by an increasing union of open quasi-projective
  \(k\)-schemes.

  \emph{Step 2: bound twists by \(m\)-regularity.}
  Given \(\phi\), fix \(m\) large enough that the \(T\)-points of \(P^\phi\) are
  representable by \(\mathcal{L}\) on \(C \times_k T\) satisfying
  \(R^i \pi_{T*} \mathcal{L}(n) = 0\) for all \(i \ge 1\), \(n \ge m\). Twisting
  by \(\mathcal{O}_C(m)\) defines an automorphism of \(\Pic^\sharp_{(C/k)\et}\)
  that carries \(P^\phi\) to \(P^{\phi_0}_0\) for \(\phi_0(n) := \phi(m + n)\).
  Hence it suffices to represent each \(P^{\phi_0}_0\) by a quasi-projective
  \(k\)-scheme.

  \emph{Step 3: factor through the Abel map.}
  Consider the Abel map \(A_{C/k} : \mathrm{Div}_{C/k} \to \Pic^\sharp_{(C/k)\et}\)
  of \cref{lem:line_bundle_quot_correspondence}. The source \(\mathrm{Div}_{C/k}\)
  is representable by an open subscheme of \(\Hilb_{C/k}\) (and hence by an
  open subscheme of the Quot scheme of \cref{thm:quot_representable}), and
  is quasi-projective over \(k\). Pulling back along
  \(P^{\phi_0}_0 \hookrightarrow \Pic^\sharp_{(C/k)\et}\) yields an open
  subscheme \(Z := P^{\phi_0}_0 \times_{\Pic^\sharp_{(C/k)\et}} \mathrm{Div}_{C/k}\)
  of \(\mathrm{Div}_{C/k}\), which is again quasi-projective. The projection
  \(\alpha : Z \to P^{\phi_0}_0\) is a surjection of \'etale sheaves:
  given a \(T\)-point \(\lambda\) of \(P^{\phi_0}_0\) represented by an \'etale
  cover \(T' \to T\) and an invertible sheaf \(\mathcal{L}'\) on \(C \times_k
  T'\) with \(H^1(C_{t'}, \mathcal{L}'_{t'}) = 0\), the projective bundle
  \(\mathbb{P}(\mathcal{Q}) \to T'\) associated to
  \(\mathcal{Q} := \pi_{T'*}\mathcal{L}'\) is smooth over \(T'\); an \'etale
  cover \(T_1 \to T'\) admits a \(T'\)-map \(T_1 \to \mathbb{P}(\mathcal{Q})\)
  (EGA IV 17.16.3 (ii)) whose composition with the universal divisor map
  \(\mathbb{P}(\mathcal{Q}) \to Z\) produces the required \(T_1\)-point of \(Z\).
  Moreover, the first projection \(Z \times_{P^{\phi_0}_0} Z \to Z\) is
  smooth and proper.

  \emph{Step 4: descend by the smooth-proper quotient lemma.}
  Apply Kleiman~\S 4, Lem.~4.9: given a surjection \(\alpha : Z \to P\)
  of \'etale sheaves with \(Z\) quasi-projective, \(R := Z \times_P Z\)
  representable, and the first projection \(R \to Z\) smooth and proper,
  \(P\) is representable by a quasi-projective scheme. In our case
  \(P = P^{\phi_0}_0\), \(Z\) and \(R\) as above, so \(P^{\phi_0}_0\) is
  quasi-projective. Reassembling over \(\phi \in \mathbb{Q}[n]\) and
  over \(m\), the disjoint union of representable \(P^\phi\) pieces represents
  \(\Pic^\sharp_{(C/k)\et}\). The total scheme is separated (as a disjoint
  union of separated quasi-projectives) and locally of finite type. This
  is Kleiman~\S 4, Cor.~4.18.3 in the \(S = \Spec k\),
  \(X = C\) complete case.
\end{proof}

\subsection{The smooth-proper quotient lemma}
\label{subsec:fga_pic_qt}

The structural lemma underlying step 4 above is recorded separately
because it is the abstract content allowing descent from
\(Z = \mathrm{Div}^{\phi_0}_0\) to the representing Picard piece.

\begin{lemma}[Smooth-proper quotient of \'etale sheaves]\leanok
  \label{lem:smooth_proper_quotient}
  \lean{AlgebraicGeometry.Scheme.PicScheme.smoothProperQuotient}
  \uses{def:has_smooth_proper_quotient}
  % SOURCE: [Kleiman], ``The Picard scheme'', \S 4, Lem.~lm:qt;
  % arXiv:math/0504020, p.~31 (read from
  % references/kleiman-picard-src/kleiman-picard.tex, L2368-L2375).
  % Source identifier: \begin{lem}\label{lm:qt}\end{lem}.
  \textit{Source: [Kleiman], ``The Picard scheme'', \S 4, Lem.~4.9.}
  Let \(\alpha : Z \to P\) be a morphism of \'etale sheaves on \((\Sch/k)\).
  Set \(R := Z \times_P Z\). Suppose:
  \begin{enumerate}
    \item \(\alpha\) is a surjection of \'etale sheaves;
    \item \(Z\) is representable by a quasi-projective \(k\)-scheme;
    \item \(R\) is representable by a \(k\)-scheme;
    \item the first projection \(R \to Z\) is representable by a smooth
      and proper map.
  \end{enumerate}
  Then \(P\) is representable by a quasi-projective \(k\)-scheme, and \(\alpha\)
  is representable by a smooth map.
\end{lemma}

\begin{proof}
  By (ii) the structure map \(Z \to \Spec k\) is quasi-projective, hence
  separated, so the first projection \(Z \times_k Z \to Z\) is separated.
  By (iv) the first projection \(R \to Z\) is proper; it factors through
  \(h : R \to Z \times_k Z\) over the diagonal embedding (the second
  component is the second projection of \(R \to Z \times_P Z\)). The map
  \(h\) is therefore proper, and a monomorphism (because \(R\) is the graph
  of an equivalence relation on \(Z\), hence injective on \(T\)-points).
  By EGA IV 8.11.5, a proper monomorphism is a closed immersion: \(h\) is a
  closed embedding.

  Hence \(R \hookrightarrow Z \times_k Z\) is a flat (by (iv)) and proper
  equivalence relation on the quasi-projective \(k\)-scheme \(Z\). By the
  theory of flat-and-proper effective equivalence relations on
  quasi-projective schemes (Altman--Kleiman; the construction places
  the graph \(R \subseteq Z \times_k Z\) inside the universal subscheme of
  the Hilbert scheme \(\Hilb_{Z/k}\) and descends to a closed subscheme
  \(Q \subseteq \Hilb_{Z/k}\) representing the quotient), there exist a
  quasi-projective \(k\)-scheme \(Q\) and a faithfully flat, projective map
  \(Z \to Q\) with \(R = Z \times_Q Z\).

  The map \(Z \to Q\) is flat and proper; its fibers, up to field
  extension, coincide with those of \(R \to Z\), which are smooth by
  hypothesis. Hence \(Z \to Q\) is smooth.

  Finally, \(Z \to Q\) is a surjection of \'etale sheaves (an \'etale
  cover of any \(T\)-point of \(Q\) pulls back to a \(T'\)-point of \(Z\), by
  EGA IV 17.16.3 (ii) applied to \(Z \times_Q T \to T\)). In the category
  of \'etale sheaves, both \(Q\) and \(P\) are coequalisers of the pair
  \(R \rightrightarrows Z\); the coequaliser is unique up to unique
  isomorphism, so \(Q\) represents \(P\). Smoothness of \(\alpha : Z \to P\)
  is the smoothness of \(Z \to Q\) just shown.
\end{proof}

\begin{remark}[Quasi-projectivity is load-bearing]
  \label{rem:smooth_proper_quotient_hypothesis}
  \uses{lem:smooth_proper_quotient}
  Hypothesis (ii) of \cref{lem:smooth_proper_quotient} --- that \(Z\) be
  representable by a \emph{quasi-projective} \(k\)-scheme --- cannot be weakened
  to bare representability. A Hironaka-type free \(\mathbb{Z}/2\)-action on a
  smooth proper non-projective threefold gives a smooth proper equivalence
  relation satisfying (i), (iii) and (iv) whose quotient \'etale sheaf is an
  algebraic space and not a scheme.

  This is what separates the two routes of \cref{sec:fga_pic_setup}. The
  quotient route must supply quasi-projectivity of the Abel-map slice; the
  Milne--Koll\'ar route quotients only by a finite Galois group, under the
  hypothesis that every finite orbit lies in an affine open, which is a
  condition the Hironaka example fails.
\end{remark}

\section{Group-scheme structure}
\label{sec:fga_pic_group_structure}

The Picard scheme \(\Pic_{C/k}\) inherits an abelian group structure from
the tensor product of invertible sheaves on \(C \times_k T\). The relative
Picard presheaf \(T \mapsto \Pic(C \times_k T)/H_T\) is already an abelian
group at every \(T\) (the quotient of an abelian group by a subgroup),
functorially in \(T\); this lifts to the \'etale sheafification, and then
to a group-scheme structure on the representing scheme by Yoneda.

\begin{theorem}[\(\Pic_{C/k}\) is a \(k\)-group scheme]\leanok
  \label{thm:pic_is_group_scheme}
  \lean{AlgebraicGeometry.Scheme.PicScheme.groupSchemeStructure}
  \uses{thm:fga_pic_representability, thm:relative_pic_quotient_well_defined,
        def:pic_scheme, def:rel_pic_sharp}

  % SOURCE: [Kleiman], ``The Picard scheme'', \S 2, Def.~df:Pfs (the relative
  % Picard functor takes values in abelian groups) and \S 4 surrounding
  % text (group-scheme structure on \(\IPic_{X/S}\) via tensor product);
  % arXiv:math/0504020, pp.~17, 26 (read from
  % references/kleiman-picard-src/kleiman-picard.tex, L1311-L1318 + L2057-L2069).
  % Source identifiers: \begin{dfn}\label{df:Pfs}\end{dfn} and
  % \begin{dfn}\label{df:Psch}\end{dfn}.
  \textit{Source: [Kleiman], ``The Picard scheme'', \S 2, Def.~2.2 +
  \S 4, Def.~4.1 (the Picard scheme).}
  Let \(\Pic_{C/k}\) be the Picard scheme of
  \cref{thm:fga_pic_representability}. Then \(\Pic_{C/k}\) is canonically an
  abelian group scheme over \(k\): the tensor product
  \([\mathcal{L}] + [\mathcal{M}] := [\mathcal{L} \otimes \mathcal{M}]\)
  and inverse \(-[\mathcal{L}] := [\mathcal{L}^{-1}]\) on
  \(\Pic^\sharp_{C/k}(T)\) (well defined because both operations
  preserve isomorphism classes and descend through the subgroup
  \(H_T = \pi_T^*\Pic(T)\) of \cref{thm:relative_pic_quotient_well_defined})
  are natural in \(T\), hence by Yoneda are represented by morphisms
  \(m : \Pic_{C/k} \times_k \Pic_{C/k} \to \Pic_{C/k}\),
  \(i : \Pic_{C/k} \to \Pic_{C/k}\), and
  \(e : \Spec k \to \Pic_{C/k}\) (the trivial line bundle class) satisfying
  the abelian-group axioms. In particular, \(\Pic_{C/k}\) is a \(k\)-group
  scheme locally of finite type.
\end{theorem}

\begin{proof}
\leanok
  The relative Picard functor \(\Pic^\sharp_{C/k}\) of
  \cref{def:rel_pic_sharp} is abelian-group valued, with sum
  \([\mathcal{L}] + [\mathcal{M}] = [\mathcal{L} \otimes \mathcal{M}]\),
  inverse \(-[\mathcal{L}] = [\mathcal{L}^{-1}]\), and unit
  \(0 = [\mathcal{O}_{C \times_k T}]\); these operations descend through
  the quotient by \(H_T = \pi_T^*\Pic(T)\) because \(\pi_T^*\) is itself a
  group homomorphism (\cref{thm:rel_pic_addcommgroup_via_tensorobj}).

  By \cref{thm:fga_pic_representability}, \(\Pic^\sharp_{C/k}\) is
  represented by the scheme \(\Pic_{C/k}\). The Yoneda embedding is fully
  faithful and preserves limits, so a representing object of a
  group-valued presheaf is canonically a group object: the natural
  transformations
  $\Pic^\sharp_{C/k} \times \Pic^\sharp_{C/k} \to \Pic^\sharp_{C/k}$
  (multiplication), \(\Pic^\sharp_{C/k} \to \Pic^\sharp_{C/k}\) (inverse),
  and the unit map \(h_{\Spec k} \to \Pic^\sharp_{C/k}\) pick out unique
  morphisms of representing schemes \(m\), \(i\), \(e\) as displayed, and
  the abelian-group axioms hold at the level of \(T\)-points functorially
  in \(T\), hence on the representing scheme. The local-finite-type
  property is part of \cref{thm:fga_pic_representability}.
\end{proof}

\section{The Picard scheme}
\label{sec:fga_pic_named_definition}

The two preceding theorems determine the following named scheme.

\begin{definition}[The Picard scheme]\leanok
  \label{def:pic_scheme}
  \lean{AlgebraicGeometry.Scheme.PicScheme}
  \uses{def:has_pic_scheme, def:inst_has_pic_scheme}

  % SOURCE: [Kleiman], ``The Picard scheme'', \S 4, Def.~df:Psch;
  % arXiv:math/0504020, p.~26 (read from
  % references/kleiman-picard-src/kleiman-picard.tex, L2057-L2062).
  % Source identifier: \begin{dfn}\label{df:Psch}\end{dfn}.
  \textit{Source: [Kleiman], ``The Picard scheme'', \S 4, Def.~4.1.}
  Let \(C/k\) be a smooth proper geometrically integral curve over a field
  \(k\), with a \(k\)-rational point. The \emph{Picard scheme} \(\Pic_{C/k}\)
  is the \(k\)-scheme representing the relative Picard functor
  \(\Pic^\sharp_{C/k}\) (equivalently, by Kleiman \S 2 Thm.~2.5 under the
  rational point, its \'etale sheafification), supplied by
  \cref{thm:fga_pic_representability}, equipped with the abelian-group-scheme
  structure of \cref{thm:pic_is_group_scheme}. The scheme is separated,
  locally of finite type over \(k\), and a disjoint union of open
  quasi-projective \(k\)-subschemes.
\end{definition}

\section{Construction interfaces}
\label{sec:fga_pic_lean_encoding}

The constructions above fit together through the following interfaces. The
relative Picard presheaf is obtained from the quotient of invertible sheaves
on \(C\times_kT\) by pullbacks from \(T\), and its set-valued version is used
when applying representability. The relative-divisor functor is the
specialisation of the effective Cartier-divisor functor to
\(C\to\operatorname{Spec}k\). A divisor family \(\langle\mathcal F,q\rangle\)
is sent by the Abel map to the inverse class of its kernel:
\[
  \langle\mathcal F,q\rangle\longmapsto-[\ker q].
\]
The FGA theorem applies Hilbert-polynomial stratification, a regularity bound,
the Abel-map factorisation, and the smooth-proper quotient lemma. Tensor
product and dualisation of line bundles then induce addition and inversion on
the representing scheme.

\begin{itemize}
  \item The Picard scheme \cref{def:pic_scheme} is selected by the
    existence assertion \cref{def:has_pic_scheme}.
  \item The functor \cref{def:pic_sharp_carrier} is the set-valued form of
    the relative Picard presheaf of \cref{chap:Picard_RelPicFunctor}.
  \item The Abel map of \cref{lem:line_bundle_quot_correspondence} is the natural
    transformation \(\mathrm{Div}_{C/k} \to \Pic^\sharp_{C/k}\) sending a
    relative divisor family \(\langle \mathcal F, q\rangle\) to
    \(-[\ker q]\), the class of its associated invertible sheaf.
  \item The Altman--Kleiman lemma \cref{lem:smooth_proper_quotient} gives
    descent of a flat-and-proper equivalence relation \(R \subseteq Z
    \times_k Z\) on a quasi-projective \(k\)-scheme \(Z\) to a
    quasi-projective quotient \(Q\) representing the coequalising sheaf.
  \item The theorem \cref{thm:fga_pic_representability} supplies the
    representing scheme.
  \item Yoneda transport of the abelian-group structure on
    \(\Pic^\sharp_{C/k}\) gives the group-scheme structure of
    \cref{thm:pic_is_group_scheme}.
\end{itemize}

The effective-quotient theorem for flat and proper equivalence relations on
quasi-projective schemes is recorded as
\cref{lem:smooth_proper_quotient}. All other interfaces are supplied by the
relative Picard, divisor, and Quot constructions.

\section{Auxiliary existence predicates}
\label{sec:fga_pic_typeclass_scaffolding}

The public constructions use a small collection of existence predicates. Each
predicate records one mathematical object or property needed by the preceding
representability argument, and the associated definitions select the object
when the predicate is inhabited.

\subsection{Basic functors}
\label{subsec:fga_pic_functor_carriers}

\begin{definition}[\(k\)-rational point]\leanok
  \label{def:has_rational_point}
  \lean{AlgebraicGeometry.Scheme.HasRationalPoint}
  The predicate \(\mathrm{HasRationalPoint}(C)\) asserts that the structural
  morphism \(C \to \Spec k\) admits a section
  \(\sigma : \Spec k \to C\), i.e.\ that \(C\) has a \(k\)-rational point.
  This is the standing hypothesis of the chapter (see
  \cref{sec:fga_pic_setup}): by Kleiman \S 2, Thm.~2.5, it makes the plain
  relative Picard functor \(\Pic^\sharp_{C/k}\) coincide with its \'etale
  (and fppf) sheafification, so representability can be stated against
  \cref{def:rel_pic_sharp} directly. For a pointed curve \((C,P)\), the
  pointing \(P\) supplies this section.
\end{definition}

\begin{proof}\leanok
  This is a definition.
\end{proof}

\begin{definition}[Set-valued relative Picard functor]\leanok
  \label{def:pic_sharp_carrier}
  \lean{AlgebraicGeometry.Scheme.PicScheme.picSharp}
  \uses{def:rel_pic_sharp}
  The set-valued shadow \(\Pic^\sharp_{C/k} : (\Sch/k)^{op} \to \Set\) of the
  group-valued relative Picard presheaf of \cref{def:rel_pic_sharp}
  (composition with the forgetful functor \(\Ab \to \Set\)).

  A representing object therefore carries natural bijections
  \[
    (T \to \Pic_{C/k}) \simeq
    \Pic(C \times_k T)/\pi_T^*\Pic(T).
  \]
\end{definition}

\begin{proof}\leanok
  \uses{def:rel_pic_sharp}
  Immediate from \cref{def:rel_pic_sharp}: apply the forgetful functor to
  the group-valued presheaf.
\end{proof}

\begin{definition}[Existence of the relative-divisor functor]\leanok
  \label{def:has_div_functor}
  \lean{AlgebraicGeometry.Scheme.PicScheme.HasDivFunctor}
  The predicate \(\mathrm{HasDivFunctor}(C)\) asserts that the category of
  presheaves of types on \((\Sch/k)^{op}\) contains the
  relative-effective-divisor functor \(\mathrm{Div}_{C/k}\), sending a
  \(k\)-scheme \(T\) to the set of relative effective Cartier divisors on
  \(C \times_k T\) flat over \(T\). It is the existence assertion for
  \(\mathrm{Div}_{C/k}\).
\end{definition}

\begin{proof}
  This is the defining predicate.
\end{proof}

\begin{definition}[The relative-divisor functor]\leanok
  \label{def:div_functor_carrier}
  \lean{AlgebraicGeometry.Scheme.PicScheme.divFunctor}
  \uses{def:div_functor}
  The presheaf \(\mathrm{Div}_{C/k} : (\Sch/k)^{op} \to \Set\) is the
  relative-divisor functor \(\Div_{C/k}\) of \cref{def:div_functor},
  instantiated at the structural morphism \(C \to \Spec k\): it sends a
  \(k\)-scheme \(T\) to the set of relative effective Cartier divisors on
  \(C \times_k T\) flat over \(T\).
\end{definition}

\begin{proof}\leanok
  \uses{def:div_functor}
  Immediate from \cref{def:div_functor}: instantiate at
  \(X = C\), \(S = \Spec k\), \(\pi\) the structural morphism.
\end{proof}

\begin{definition}[Relative-divisor functor exists]\leanok
  \label{def:inst_has_div_functor}
  \lean{AlgebraicGeometry.Scheme.PicScheme.instHasDivFunctor}
  \uses{def:has_div_functor, def:div_functor_carrier}
  The relative-divisor functor of \cref{def:div_functor_carrier} witnesses
  \cref{def:has_div_functor} for every such curve \(C/k\).
\end{definition}

\begin{proof}\leanok
  \uses{def:div_functor_carrier}
  The functor \(\mathrm{Div}_{C/k}\) of \cref{def:div_functor_carrier} is an
  object of the presheaf category, witnessing the existential.
\end{proof}

\subsection{Abel-map predicate}
\label{subsec:fga_pic_abel_map_carrier}

\begin{definition}[Abel-map data]\leanok
  \label{def:has_abel_map}
  \lean{AlgebraicGeometry.Scheme.PicScheme.HasAbelMap,
    AlgebraicGeometry.Scheme.PicScheme.abelMap}
  \uses{def:div_functor_carrier, def:pic_sharp_carrier}
  The predicate \(\mathrm{HasAbelMap}(C)\) carries a natural
  transformation of presheaves \(\mathrm{abel} \colon \mathrm{Div}_{C/k}
  \to \Pic^\sharp_{C/k}\) between the relative-divisor functor of
  \cref{def:div_functor_carrier} and the relative Picard functor of
  \cref{def:pic_sharp_carrier}; \(\mathrm{abelMap} :=
  \mathrm{HasAbelMap.abel}\) denotes this transformation.
\end{definition}

\begin{proof}
  This is a definition.
\end{proof}

\begin{definition}[Existence of the Abel map]\leanok
  \label{def:inst_has_abel_map}
  \lean{AlgebraicGeometry.Scheme.PicScheme.instHasAbelMap}
  \uses{def:has_abel_map, def:abel_map_witness}
  For every such curve \(C/k\), the transformation
  \(A^{\mathrm{raw}}_{C/k}\) from \cref{def:abel_map_witness} supplies
  \cref{def:has_abel_map}.
\end{definition}

\begin{proof}
\leanok
  \uses{def:abel_map_witness}
  The natural transformation constructed in \cref{def:abel_map_witness} is a
  term of the required type
  \(\mathrm{Div}_{C/k} \to \Pic^\sharp_{C/k}\), and hence supplies the data of
  \cref{def:has_abel_map}.
\end{proof}

\subsection{Smooth-proper quotient predicate}
\label{subsec:fga_pic_quotient_carrier}

\begin{definition}[Representability of the smooth-proper quotient]\leanok
  \label{def:has_smooth_proper_quotient}
  \lean{AlgebraicGeometry.Scheme.PicScheme.HasSmoothProperQuotient}
  For a morphism \(\alpha : Z \to P\) of presheaves of types on
  \((\Sch/k)^{op}\), the predicate \(\mathrm{HasSmoothProperQuotient}(\alpha)\)
  asserts that the target \(P\) is representable. It records the conclusion
  of the Altman--Kleiman lemma \cref{lem:smooth_proper_quotient} for the
  quotient occurring in the FGA construction.
\end{definition}

\begin{proof}\leanok
  This is a definition.
\end{proof}

\subsection{Picard-scheme predicates}
\label{subsec:fga_pic_scheme_carriers}

\begin{definition}[Existence of the Picard scheme]\leanok
  \label{def:has_pic_scheme}
  \lean{AlgebraicGeometry.Scheme.HasPicScheme}
  \uses{def:pic_sharp_carrier}
  The predicate \(\mathrm{HasPicScheme}(C)\) asserts that there exists a \(k\)-scheme
  \(X\), \emph{locally of finite type over \(k\)}, carrying a
  representability witness for the relative Picard functor
  \(\Pic^\sharp_{C/k}\) of \cref{def:pic_sharp_carrier}; in other words,
  that the Picard scheme \(\Pic_{C/k}\) exists and is locally of finite
  type. It is the existence assertion from which \cref{def:pic_scheme}
  extracts the representing object. The local-finiteness clause belongs to
  the same existence package as representability: Kleiman \S 4,
  Thm.~th:main exhibits \(\Pic_{C/k}\) as a disjoint union of open
  quasi-projective \(k\)-subschemes. The witness
  \cref{def:inst_has_pic_scheme} is conditional on the \(k\)-rational point
  of \cref{def:has_rational_point}.
\end{definition}

\begin{proof}
  This is the defining existence predicate.
\end{proof}

\begin{definition}[Picard scheme exists]\leanok
  \label{def:inst_has_pic_scheme}
  \lean{AlgebraicGeometry.Scheme.instHasPicScheme}
  \uses{def:has_pic_scheme, def:has_rational_point}
  Let \(C/k\) be a smooth proper geometrically integral curve with a
  \(k\)-rational point. Then the relative Picard functor of \(C\) is
  represented by a scheme locally of finite type over \(k\).
\end{definition}

\begin{proof}
  Apply the FGA representability construction of
  \cref{thm:fga_pic_representability}. The rational point rigidifies line
  bundles along the section, and the bounded Hilbert-polynomial strata glue to
  the required locally finite type representing scheme.
\end{proof}

\begin{definition}[Representability of \(\Pic^\sharp_{C/k}\)]\leanok
  \label{def:pic_sharp_representable}
  \lean{AlgebraicGeometry.Scheme.PicScheme.PicSharpRepresentable}
  \uses{def:pic_sharp_carrier, def:pic_scheme}
  The predicate \(\mathrm{PicSharpRepresentable}(C)\) asserts that the specific scheme
  \(\Pic_{C/k}\) of \cref{def:pic_scheme} is a representing object for
  \(\Pic^\sharp_{C/k}\). It is the assertion from which the representability
  witness \cref{thm:fga_pic_representability} is extracted.
\end{definition}

\begin{proof}
  This is a definition.
\end{proof}

\begin{definition}[Representability identification]\leanok
  \label{def:inst_pic_sharp_representable}
  \lean{AlgebraicGeometry.Scheme.PicScheme.instPicSharpRepresentable}
  \uses{def:pic_sharp_representable, def:has_pic_scheme}
  The scheme \(\Pic_{C/k}\) is the chosen witness of
  \cref{def:has_pic_scheme}; the first component of that witness identifies
  it as a representing object for \(\Pic^\sharp_{C/k}\).
\end{definition}

\begin{proof}\leanok
  \uses{def:has_pic_scheme}
  This is the representing component of the existence statement in
  \cref{def:has_pic_scheme}.
\end{proof}

\begin{definition}[Local finite type of \(\Pic_{C/k}\)]\leanok
  \label{def:pic_scheme_lft}
  \lean{AlgebraicGeometry.Scheme.PicScheme.PicSchemeLocallyOfFiniteType}
  \uses{def:pic_scheme, def:has_pic_scheme}
  The predicate
  \[
    \mathrm{PicSchemeLocallyOfFiniteType}(C)
  \]
  asserts that the
  structural morphism of the Picard scheme \(\Pic_{C/k}\) of
  \cref{def:pic_scheme} is locally of finite type --- part (1) of Kleiman
  \S 4, Thm.~th:main (\(\Pic_{C/k}\) is a disjoint union of open
  quasi-projective \(k\)-subschemes). It is the hypothesis the
  identity-component construction of
  \cref{chap:Picard_IdentityComponent} uses to specialise
  \(G^0\) to \(G = \Pic_{C/k}\). It follows from the local-finiteness
  component of \cref{def:has_pic_scheme}.
\end{definition}

\begin{proof}\leanok
  This is a definition.
\end{proof}

\begin{definition}[Local-finiteness assertion]\leanok
  \label{def:inst_pic_scheme_lft}
  \lean{AlgebraicGeometry.Scheme.PicScheme.instPicSchemeLocallyOfFiniteType}
  \uses{def:pic_scheme_lft, def:has_pic_scheme}
  Whenever \cref{def:has_pic_scheme} holds, the second component of its
  chosen witness gives local finiteness of \(\Pic_{C/k}\). Thus
  representability and local finiteness are supplied by one existence
  package, as in Kleiman \S 4, Thm.~th:main.
\end{definition}

\begin{proof}\leanok
  \uses{def:has_pic_scheme}
  This is the local-finiteness component of the existence statement in
  \cref{def:has_pic_scheme}.
\end{proof}

\section{Dependencies of the representability theorem}
\label{sec:fga_pic_sorry_closure_order}

The construction of \(\Pic_{C/k}\) separates into three mathematical inputs:
the relative-divisor functor, its Abel map to the relative Picard functor, and
the FGA existence theorem. The first two supply concrete functors and a
natural transformation. The third combines them with regularity, cohomology
and base change, and effective descent to produce a scheme locally of finite
type representing \(\Pic^\sharp_{C/k}\).

\subsection{The Picard-scheme existence theorem}
\label{subsec:sorry_has_pic_scheme}

For a
smooth proper geometrically integral curve \(C/k\) with a \(k\)-rational
point, the relative Picard functor \(\Pic^\sharp_{C/k}\)
(\cref{def:pic_sharp_carrier}) is representable by a \(k\)-scheme. The
proof first applies Kleiman \S 2, Thm.~2.5, using the section to identify the
plain functor with its \'etale sheafification, and then uses Kleiman's four
steps
(\cref{thm:fga_pic_representability}): Hilbert-polynomial stratification,
\(m\)-regularity bound, Abel-map factorisation through
\(\mathrm{Div}_{C/k}\), and the smooth-proper quotient
(\cref{lem:smooth_proper_quotient}). Its inputs are the Quot/Hilbert
construction of \cref{chap:Picard_QuotScheme}, representability of
\(\mathrm{Div}_{C/k}\) as in \cref{rem:div_functor_representability}, the
regularity result \cref{cor:flattening_stratification_curve}, the base-change
result \cref{lem:cech_flat_base_change}, and Altman--Kleiman effective descent.
The resulting representing object gives
natural bijections
\[
  (T\to\Pic_{C/k})\simeq
  \Pic(C\times_kT)/\pi_T^*\Pic(T),
\]
including the dual-number case \(T=\Spec k[\epsilon]\) used in the
tangent-space calculation.

\subsection{The relative-divisor input}
\label{subsec:sorry_has_div_functor}

The relative-divisor functor \(\Div_{X/S}\) of
\cref{def:div_functor} is the subfunctor of the Hilbert functor
\(\Hilb_{C/k} = \Quot_{\mathcal{O}_C/C/k}\) of relative effective Cartier
divisors, in the invertible-kernel quotient encoding of
\cref{def:div_family}. The presheaf \(\mathrm{Div}_{C/k}\) of
\cref{def:div_functor_carrier} is its instantiation at
\(C\to\Spec k\). Its divisor condition is stable under base change by
\cref{lem:relative_divisor_base_change}, while representability of
\(\Div_{C/k}\) by an open subscheme of the Hilbert scheme (Kleiman \S 3,
Thm.~th:repDiv) is the content of
\cref{rem:div_functor_representability}.

\subsection{The Abel-map input}
\label{subsec:sorry_has_abel_map}

The class \cref{def:has_abel_map} carries the natural transformation whose
value on a divisor family is
\[
  \mathrm{abelMapWitness}(C)\langle\mathcal F,q\rangle=-[\ker q].
\]
It is the composite of the kernel-class natural transformation
\(\mathrm{abelKernelNatTrans}(C)\) (naturality via the kernel--flat-base-change
comparison isomorphism) with the group-presheaf negation
\(\mathrm{picNeg}(C)\). Thus the Abel map has the defining property in
\cref{lem:line_bundle_quot_correspondence}; it is a natural transformation
of the presheaves in \cref{def:div_functor_carrier} and
\cref{def:pic_sharp_carrier}.

\subsection{Assembly of the inputs}
\label{subsec:fga_pic_closure_order_summary}

The FGA construction consumes the relative-divisor functor and Abel map,
then applies the regularity, base-change, and effective-quotient results of
\cref{thm:fga_pic_representability}. A pointing
\(P:\mathbf 1\to C\) supplies the rational-point hypothesis used to identify
the plain relative Picard functor with its sheafification. The resulting
representability and local-finiteness statements feed the group-scheme and
Jacobian constructions.

\section{Related constructions}
\label{sec:fga_pic_out_of_scope}

The representable Picard scheme and its group structure interact with the
following constructions, treated in the indicated companion chapters.
\begin{itemize}
  \item \textbf{Identity component \(\Pic^0_{C/k}\)} --- the open and
    closed subgroup scheme of \(\Pic_{C/k}\) containing the unit and
    consisting of those classes lying in the connected component of
    \(\Pic_{C/k}\). Its dimension, smoothness, and abelian-variety
    structure are treated in \cref{chap:Picard_IdentityComponent} and
    \cref{chap:Picard_Pic0AbelianVariety}.
  \item \textbf{Quot scheme representability} --- consumed by
    \cref{thm:fga_pic_representability} via \texttt{thm:quot\_representable}
    (\cref{chap:Picard_QuotScheme}), not re-proved here.
  \item \textbf{\'Etale sheafification.} The relative Picard presheaf and its
    topology-parametric sheafification are constructed in
    \cref{chap:Picard_RelPicFunctor}. Under the
    standing rational-point hypothesis, Kleiman \S 2, Thm.~2.5 identifies
    the plain relative functor with the \'etale sheafification.
  \item \textbf{Universal Poincar\'e sheaf} (Kleiman~\S 4, Ex.~4.3)
    --- exists when \(\Pic_{C/k}\) represents \(\Pic_{C/k}\) itself and
    \(C/k\) has a section. Under the chapter's standing rational-point
    hypothesis both conditions in fact hold (Thm.~2.5 identifies the
    plain functor with the represented sheaf), so the Poincar\'e sheaf
    is then available as an associated construction.
  \item \textbf{Mumford's example} (Kleiman~\S 4, Eg.~4.14), counter-examples
    to representability for non-integral geometric fibers --- not
    relevant to a smooth proper geometrically integral curve.
\end{itemize}

\section{Dependency structure}
\label{sec:fga_pic_consistency_check}

The principal constructions form a connected dependency chain rooted in
\cref{chap:Picard_LineBundlePullback}, \cref{chap:Picard_QuotScheme}, and
\cref{chap:Picard_RelPicFunctor}:
\begin{itemize}
  \item \cref{lem:line_bundle_quot_correspondence} uses
    \cref{def:line_bundle_on_product} and
    \cref{thm:pullback_natural}.
  \item \cref{lem:smooth_proper_quotient} is the structural quotient
    statement used by the representability argument.
  \item \cref{thm:fga_pic_representability} uses
    \cref{lem:line_bundle_quot_correspondence},
    \texttt{thm:quot\_representable},
    \texttt{def:rel\_pic\_sharp},
    \cref{def:has_rational_point},
    \cref{thm:relative_pic_quotient_well_defined}, and
    \cref{lem:smooth_proper_quotient}.
  \item \cref{thm:pic_is_group_scheme} uses
    \cref{def:pic_scheme}, \cref{thm:fga_pic_representability},
    \texttt{def:rel\_pic\_sharp}, and
    \cref{thm:relative_pic_quotient_well_defined}.
  \item \cref{def:pic_scheme} uses
    \cref{def:has_pic_scheme} and \cref{def:inst_has_pic_scheme}.
\end{itemize}
The Quot-scheme theorem used above is proved in
\cref{chap:Picard_QuotScheme}, while the relative Picard functor is defined in
\cref{chap:Picard_RelPicFunctor}.
